\documentclass[11pt]{article}
% =====================================================================
%  JMLR companion theory paper -- DRAFT (Track T; T2-4 assembly in progress)
%  Template-agnostic article; port to jmlr2e.sty later.
%
%  Ported from ../theory_notes.tex (Rev.2, 3rd-pass reviewed). Convention:
%  standard measure-/Markov-theoretic facts are STATED WITH CITATION (not
%  reproved) at paper altitude; the substantive/novel results carry full
%  proofs. Notation unified to bold theta/u throughout (theory_notes mixes
%  bold in secs 1-8 with plain in the contraction/steering secs).
%
%  Static-check convention (no LaTeX toolchain in-session):
%   balanced \begin/\end; every \ref has a \label; every \cite has a \bibitem.
% =====================================================================

\usepackage[margin=1in]{geometry}
\usepackage{amsmath,amssymb,amsthm,mathtools}
\usepackage{bm}
\usepackage{booktabs}
\usepackage{array}
\usepackage[round,authoryear]{natbib}
\usepackage{hyperref}

\DeclareMathOperator{\intr}{int}

\theoremstyle{plain}
\newtheorem{theorem}{Theorem}
\newtheorem{lemma}[theorem]{Lemma}
\newtheorem{proposition}[theorem]{Proposition}
\newtheorem{corollary}[theorem]{Corollary}
\newtheorem{fact}[theorem]{Fact}
\theoremstyle{definition}
\newtheorem{definition}{Definition}
\newtheorem{assumption}{Assumption}
\theoremstyle{remark}
\newtheorem{remark}{Remark}

\newcommand{\R}{\mathbb{R}}
\newcommand{\N}{\mathbb{N}}
\newcommand{\E}{\mathbb{E}}
\newcommand{\Prob}{\mathbb{P}}
\newcommand{\cF}{\mathcal{F}}
\newcommand{\cP}{\mathcal{P}}
\newcommand{\Om}{\Omega}
\newcommand{\Leb}{\mathrm{Leb}}
\newcommand{\TV}{\mathrm{TV}}
\newcommand{\Normal}{\mathcal{N}}
\newcommand{\xobs}{\bm{x}_{\mathrm{obs}}}
\newcommand{\bth}{\bm{\theta}}
\newcommand{\bu}{\bm{u}}
\newcommand{\bz}{\bm{z}}
\newcommand{\push}{{}_{\#}}
\newcommand{\norm}[1]{\lVert #1 \rVert}
\newcommand{\abs}[1]{\lvert #1 \rvert}
\newcommand{\einc}{\varepsilon_{\mathrm{inc}}}
\newcommand{\edist}{\varepsilon_{\mathrm{dist}}}

% ---- Working title (finalize in T2-5) --------------------------------
\title{\textbf{Convergence, Consistency, and a Computable Compatibility\\
Certificate for Amortized Pseudo-Gibbs Samplers with Learned Conditionals}}
\author{Anonymous (companion to the A-DCVAE application paper)}
\date{Draft --- \today}

\begin{document}
\maketitle

\begin{abstract}
We study amortized \emph{pseudo-Gibbs} samplers built by alternating two \emph{separately learned}
neural conditional samplers $Q_H(\bu\mid\bth)$ and $Q_P(\bth\mid\bu)$ on a \emph{continuous} state
space. Unlike a true Gibbs sampler, the two learned conditionals need not be the conditionals of any
common joint: they may be \emph{incompatible}. We give a unified, non-asymptotic analysis of the
resulting chain: (i) \emph{existence and uniform ergodicity} of a unique stationary law in total
variation, from injected noise and compactness alone (no contraction assumed, Theorem~\ref{thm:A});
(ii) a \emph{Wasserstein mixing rate} governed by the condition-Lipschitz constants
(Theorem~\ref{thm:rate}); and (iii) a \emph{steering / consistency} bound on the distance from the
stationary law to the true posterior (Theorem~\ref{thm:B}). Our central object is a computable
\emph{self-consistency certificate} $\einc$, estimable from a single run of the trained model without
access to the truth, which we show (a) vanishes iff the learned conditionals are compatible
(Proposition~\ref{prop:inc}), and (b) is bounded by the individual approximation defects with an
explicit constant (Theorem~\ref{thm:incbound}). We position the results against approximate/perturbed
MCMC and dependency networks, and validate them on controlled examples.
\end{abstract}

\tableofcontents

% =====================================================================
\section{Introduction}
\label{sec:intro}

A true Gibbs sampler alternates the conditionals $p(\bu\mid\bth)$ and $p(\bth\mid\bu)$ of a
\emph{single} joint law; compatibility is automatic and the chain is exact for that joint. We study
what happens when the two conditionals are instead \emph{learned separately}, as neural samplers
$Q_H(\bu\mid\bth)$ and $Q_P(\bth\mid\bu)$, and alternated on a \emph{continuous} state space. Nothing
forces $Q_H,Q_P$ to be the conditionals of any common joint---they may be \emph{incompatible}---so
four questions arise: does the sweep converge (existence), how fast (rate), to \emph{what} and how far
from the true posterior (steering), and can one \emph{detect} incompatibility without knowing the
truth (certification)?

The continuous setting is where the difficulty concentrates. For a \emph{finite} state space a
strictly positive transition matrix already yields a unique stationary law and geometric convergence
by Perron--Frobenius; this is the regime of dependency networks \citep{heckerman2000}, whose analysis
of separately learned conditionals is discrete and asymptotic. On a continuous space positivity is not
enough: we must supply a minorization to obtain existence (Section~\ref{sec:existence}), and a
Wasserstein contraction---not total variation, which is blind to how close distinct points are---to
obtain a dimension-robust rate and a stability constant (Sections~\ref{sec:rate}--\ref{sec:steering}).
Our steering bound has the shape of a classical perturbed-MCMC bound
\citep{mitrophanov2005,rudolf2018}, and we are explicit that the machinery is standard; what is new is
the \emph{object}---separately learned, possibly incompatible conditionals in continuous space---and a
\emph{computable} certificate of self-consistency. Section~\ref{sec:positioning} makes the comparison
precise.

\paragraph{Contributions.}
\begin{enumerate}
  \item \textbf{A new object made rigorous.} A non-asymptotic, continuous-state theory for Gibbs
        sweeps assembled from \emph{separately learned, possibly incompatible} conditionals:
        existence/uniform ergodicity (Theorem~\ref{thm:A}), a Wasserstein rate (Theorem~\ref{thm:rate}),
        and a steering bound to the true posterior (Theorem~\ref{thm:B}).
  \item \textbf{A computable compatibility certificate $\einc$.} Estimable from one run of the trained
        model, with no access to the truth; it vanishes exactly when the two conditionals are
        compatible (Proposition~\ref{prop:inc}), and is dominated by the individual approximation
        defects with an \emph{explicit} constant (Theorem~\ref{thm:incbound}).
  \item \textbf{A unified frame + noise taxonomy} dissolving the apparent ``more contraction
        $\Rightarrow$ more collapse'' tension (Remark~\ref{rem:noises}).
\end{enumerate}
Simulation-based / amortized inference \citep{papamakarios2016,cranmer2020} motivates the learned
conditionals; our analysis is agnostic to how they are trained.

% =====================================================================
\section{Preliminaries and notation}
\label{sec:prelim}

We work on Borel subsets of Euclidean space. For $S\subseteq\R^d$, $\intr S$ denotes interior,
$\partial S$ boundary, $\Leb$ Lebesgue measure, and $\cP(\Om)$ the probability measures on
$(\Om,\cF)$ with $\cF=\mathcal B(\Om)$. Pushforward is $(f\push\mu)(A)=\mu(f^{-1}(A))$.
Table~\ref{tab:notation} collects the symbols that recur; we draw attention to the several
$\varepsilon$'s, which name \emph{distinct} quantities.

\begin{table}[h]
\centering
\renewcommand{\arraystretch}{1.15}
\begin{tabular}{@{}l l l@{}}
\toprule
Symbol & Meaning & Introduced \\
\midrule
\multicolumn{3}{@{}l}{\emph{Constants and defects}}\\
$\sigma_H,\sigma_P>0$ & Gaussian injection-noise scales (noise class N4) & Def.~\ref{def:condmaps}\\
$L_H,L_P$ & condition-Lipschitz constants of the \emph{learned} decoders & Asm.~\ref{ass:lip}\\
$\kappa=L_HL_P$ & sweep $W_2$-contraction constant & Thm.~\ref{thm:rate}\\
$L_H^\dagger,L_P^\dagger$ & condition-$W_2$-Lipschitz constants of the \emph{true} conditionals & Asm.~\ref{ass:truereg}\\
$\delta_H,\delta_P$ & uniform Lebesgue-density lower bounds of $Q_H,Q_P$ & Eq.~\eqref{eq:deltaH}\\
$\varepsilon=\delta_H\delta_P\Leb(\Om)$ & Doeblin minorization constant; TV rate $(1-\varepsilon)^n$ & Lem.~\ref{lem:minor}\\
$\varepsilon_H,\varepsilon_P$ & approximation defects $\sup W_2(\text{learned},\text{true})$ & Eq.~\eqref{eq:defects}\\
$\einc$ & self-consistency defect $W_2(\Pi^\rightarrow,\Pi^\leftarrow)$ (\emph{computable}) & Eq.~\eqref{eq:einc}\\
$\edist$ & distortion defect ($\sigma$-smoothing $+$ projection), folded into $\varepsilon_H,\varepsilon_P$ & Rem.~\ref{rem:edist}\\
$C,C_H,C_P$ & constants in $\einc\le C(\varepsilon_H+\varepsilon_P)$ & Thm.~\ref{thm:incbound}\\
\midrule
\multicolumn{3}{@{}l}{\emph{Kernels}}\\
$Q_H(\cdot\mid\bth),\,Q_P(\cdot\mid\bu)$ & learned conditionals & Def.~\ref{def:condmaps}\\
$p^\dagger_H(\cdot\mid\bth),\,p^\dagger_P(\cdot\mid\bu)$ & true conditionals & Asm.~\ref{ass:suptrue}\\
$T$ & full one-sweep kernel on $\Om=\Om_H\times\Om_\Theta$ & Def.~\ref{def:T}\\
$K,\,K^\dagger$ & $\bth$-marginal sweep kernels: learned, true & Thm.~\ref{thm:rate},\,\S\ref{sec:steering}\\
\midrule
\multicolumn{3}{@{}l}{\emph{Measures}}\\
$\pi^\ast$ & joint invariant law of $T$ (TV) & Thm.~\ref{thm:A}\\
$\nu^\ast$ & $\bth$-marginal invariant (the sampler's output law) & Thm.~\ref{thm:rate}\\
$\nu^\dagger=p(\bth\mid\xobs)$ & true $\bth$-posterior (the estimand) & \S\ref{sec:steering}\\
$\mu^\ast,\,\mu^\dagger$ & hidden-state marginals: forward-sweep, true & \S\ref{sec:steering}\\
$\Pi^\dagger=p(\bth,\bu\mid\xobs)$ & true joint posterior & Asm.~\ref{ass:suptrue}\\
$\Pi^\rightarrow=\nu^\ast\!\otimes Q_H,\ \Pi^\leftarrow=Q_P\!\otimes\mu^\ast$ & forward/backward assembled joints & Eq.~\eqref{eq:einc}\\
\bottomrule
\end{tabular}
\caption{Notation. The five defect symbols $\varepsilon,\varepsilon_H,\varepsilon_P,\einc,\edist$
are \emph{distinct}: $\varepsilon$ is the TV minorization constant; $\varepsilon_H,\varepsilon_P$
measure learned-vs-true conditional error; $\einc$ is the computable self-consistency defect; $\edist$
is the $\sigma$/projection distortion folded into $\varepsilon_H,\varepsilon_P$.}
\label{tab:notation}
\end{table}

\paragraph{Metrics and standard facts (cited).}
The total variation distance is $\norm{\mu-\nu}_\TV=\sup_{A}\abs{\mu(A)-\nu(A)}=\tfrac12\abs{\mu-\nu}(\Om)$;
on compact $\Om$ the quadratic Wasserstein distance is
$W_2(\mu,\nu)=(\inf_{\gamma\in\Gamma(\mu,\nu)}\int\norm{x-y}^2\,d\gamma)^{1/2}$ over couplings
$\Gamma(\mu,\nu)$. We use the following classical facts without reproof.
\begin{fact}[Standard analytic toolkit]
\label{fact:std}
\emph{(a)} $(\cP(\Om),\norm{\cdot}_\TV)$ is complete \citep{folland}; the measure action
$\mu\mapsto\mu T$ of any Markov kernel is TV-nonexpansive, and for any coupling $(X,Y)$ of $\mu,\nu$,
$\norm{\mu-\nu}_\TV\le\Prob(X\neq Y)$. \emph{(b)} Ionescu--Tulcea builds a Markov chain from any kernel
on any measurable space, with no topological assumption \citep{kallenberg}. \emph{(c)} For nonempty
compact convex $K\subset\R^d$, the Euclidean projection $\mathcal B_K$ is single-valued, $1$-Lipschitz,
fixes $K$, sends $\R^d\setminus K$ into $\partial K$, and satisfies $\mathcal B_K^{-1}(A)=A$ for
$A\subseteq\intr K$ \citep{hiriart}; moreover $\Leb(\partial K)=0$. \emph{(d)} On compact $\Om$,
$(\cP(\Om),W_2)$ is a complete metric space; $L$-Lipschitz maps satisfy
$W_2(f\push\mu,f\push\nu)\le L\,W_2(\mu,\nu)$, so coordinate marginalization is $1$-Lipschitz
\citep{villani}. \emph{(e)} \textup{(Measurable coupling selection.)} For Markov kernels $K,K'$ from
$(\mathsf S,\mathcal S)$ to a compact $\Om$ there is a measurable family $s\mapsto\gamma_s$ of optimal
couplings of $K(s,\cdot),K'(s,\cdot)$ \citep[Cor.~5.22]{villani}; gluing against any
$\rho\in\cP(\mathsf S)$ gives
$W_2(\rho K,\rho K')^2\le\int W_2(K(s,\cdot),K'(s,\cdot))^2\,\rho(ds)$. \emph{(f)}
\textup{(Banach in $W_2$.)} If $\mu\mapsto\mu K$ is a $\kappa$-contraction with $\kappa\in[0,1)$, then
$K$ has a unique invariant $\nu^\ast$ with $W_2(\mu_0 K^n,\nu^\ast)\le\kappa^n W_2(\mu_0,\nu^\ast)$.
\end{fact}

% =====================================================================
\section{The amortized pseudo-Gibbs kernel}
\label{sec:kernel}

The sparse observation $\xobs$ is fixed throughout and suppressed. The hidden-state trajectory is
$\bu\in\R^{d_H}$ and the parameter is $\bth\in\R^{d_\Theta}$, with joint state $\bz=(\bu,\bth)$.
Decoders $g_\phi:\R^{d_\Theta}\times\R^{m_H}\to\R^{d_H}$, $g_\psi:\R^{d_H}\times\R^{m_P}\to\R^{d_\Theta}$
act on latents $\bm z_H\sim\Normal(\bm0,I_{m_H})$, $\bm z_P\sim\Normal(\bm0,I_{m_P})$;
$\mathcal B_H,\mathcal B_\Theta$ are Euclidean projections onto compact convex bodies
$\Om_H,\Om_\Theta$.

\begin{definition}[Conditional stochastic maps]
\label{def:condmaps}
With $\bm\xi_H\sim\Normal(\bm0,I_{d_H})$, $\bm\xi_P\sim\Normal(\bm0,I_{d_\Theta})$ (all latents and
injection noises mutually independent) and scales $\sigma_H,\sigma_P>0$,
\[
  Q_H(\cdot\mid\bth):=(\mathcal B_H)\push\,\mathrm{law}\big(g_\phi(\bth,\bm z_H)+\sigma_H\bm\xi_H\big),
  \quad
  Q_P(\cdot\mid\bu):=(\mathcal B_\Theta)\push\,\mathrm{law}\big(g_\psi(\bu,\bm z_P)+\sigma_P\bm\xi_P\big).
\]
The pre-projection hidden step has Lebesgue density on $\R^{d_H}$
\begin{equation}\label{eq:qH-density}
  q_H(\bm y\mid\bth)=\int_{\R^{m_H}}\Normal\big(\bm y;g_\phi(\bth,\bm z_H),\sigma_H^2 I_{d_H}\big)\,
  \Normal(\bm z_H;\bm0,I_{m_H})\,d\bm z_H,
\end{equation}
an infinite Gaussian mixture; $q_P(\cdot\mid\bu)$ is analogous.
\end{definition}

\begin{definition}[The inference kernel $T$]
\label{def:T}
On $\Om:=\Om_H\times\Om_\Theta$, one sweep draws $\bu'\sim Q_H(\cdot\mid\bth)$ then
$\bth'\sim Q_P(\cdot\mid\bu')$:
\[
  T\big((\bu,\bth),A\big)=\int_{\Om_H}Q_P\big(A_{\bu'}\mid\bu'\big)\,Q_H(d\bu'\mid\bth),
  \qquad A_{\bu'}:=\{\bth':(\bu',\bth')\in A\},
\]
depending on the current state only through $\bth$.
\end{definition}

\begin{assumption}[Compact convex state space]
\label{ass:compact}
$\Om_H,\Om_\Theta$ are nonempty, compact, convex with nonempty interior, and
$\mathcal B_H,\mathcal B_\Theta$ project onto them; thus $\Om$ is compact with $\Leb(\Om)\in(0,\infty)$
and $\Leb(\partial\Om)=0$ (Fact~\ref{fact:std}(c)).
\end{assumption}
\begin{assumption}[Continuity and nondegenerate noise]
\label{ass:cont}
For the fixed $\xobs$, $g_\phi,g_\psi$ are continuous and $\sigma_H,\sigma_P>0$.
\end{assumption}
\noindent Both hold \emph{by construction}: projection gives Assumption~\ref{ass:compact};
continuous-activation networks with positive sampling noise give Assumption~\ref{ass:cont}.

\begin{lemma}[$T$ is a Markov kernel]
\label{lem:measurable}
Under Assumptions~\ref{ass:compact}--\ref{ass:cont}, $Q_H,Q_P$ and $T$ are Markov kernels.
\end{lemma}
\begin{proof}
By dominated convergence \eqref{eq:qH-density} is continuous, hence Borel, in $(\bm y,\bth)$; writing
$Q_H(A\mid\bth)=\int\mathbf 1_A(\mathcal B_H(\bm y))\,q_H(\bm y\mid\bth)\,d\bm y$, Tonelli gives
measurability in $\bth$, and $Q_H(\cdot\mid\bth)$ is a probability measure (pushforward). Symmetrically
for $Q_P$. Section integration against a kernel preserves measurability \citep[Ch.~1]{kallenberg}, so
$\bth\mapsto\int Q_P(A_{\bu'}\mid\bu')\,Q_H(d\bu'\mid\bth)=T((\bu,\bth),A)$ is measurable and
$T((\bu,\bth),\cdot)$ is a probability measure.
\end{proof}

% =====================================================================
\section{Existence and uniform ergodicity (Theorem A)}
\label{sec:existence}

\begin{definition}[Minorization]
\label{def:minor}
$T$ satisfies \emph{minorization} with $(\varepsilon,\nu)$, $\varepsilon\in(0,1]$, $\nu\in\cP(\Om)$,
if $T(x,A)\ge\varepsilon\,\nu(A)$ for all $x,A$.
\end{definition}

\begin{theorem}[Doeblin; {\citealp[Ch.~16]{meyntweedie}}]
\label{thm:doeblin}
If a Markov kernel $T$ on $(\Om,\cF)$ with Borel diagonal satisfies minorization with
$(\varepsilon,\nu)$, it has a unique invariant $\pi\in\cP(\Om)$ and
$\sup_{x}\norm{T^n(x,\cdot)-\pi}_\TV\le(1-\varepsilon)^n$, hence
$\norm{\mu_0T^n-\pi}_\TV\le(1-\varepsilon)^n$ for all $\mu_0$.
\end{theorem}

\begin{lemma}[Uniform minorization of $T$]
\label{lem:minor}
Under Assumptions~\ref{ass:compact}--\ref{ass:cont}, $T$ satisfies minorization with
$\nu(\cdot)=\Leb(\cdot\cap\intr\Om)/\Leb(\Om)$ and some $\varepsilon\in(0,1]$.
\end{lemma}
\begin{proof}
\emph{Uniform density lower bound.} Fix $R>0$ with $p_R:=\Prob(\norm{\bm z_H}\le R)>0$. Since
$\Om_\Theta\times\{\norm{\bm z_H}\le R\}$ is compact and $g_\phi$ continuous,
$M_H:=\sup\norm{g_\phi(\bth,\bm z_H)}<\infty$ over it; let $M_H':=\sup_{\bm y\in\Om_H}\norm{\bm y}$.
For $\bm y\in\Om_H,\bth\in\Om_\Theta,\norm{\bm z_H}\le R$,
$\Normal(\bm y;g_\phi(\bth,\bm z_H),\sigma_H^2 I)\ge(2\pi\sigma_H^2)^{-d_H/2}\exp(-(M_H+M_H')^2/2\sigma_H^2)=:c_H>0$.
Restricting \eqref{eq:qH-density} to $\{\norm{\bm z_H}\le R\}$,
\begin{equation}\label{eq:deltaH}
  q_H(\bm y\mid\bth)\ge c_H\,p_R=:\delta_H>0\quad\text{uniformly},\qquad
  \text{and analogously } q_P(\cdot\mid\bu)\ge\delta_P>0.
\end{equation}
\emph{Projection preserves the interior bound.} For $A\subseteq\intr\Om_H$, Fact~\ref{fact:std}(c)
gives $\mathcal B_H^{-1}(A)=A$, so $Q_H(A\mid\bth)=\int_A q_H\ge\delta_H\Leb(A)$; likewise
$Q_P(B\mid\bu)\ge\delta_P\Leb(B)$ for $B\subseteq\intr\Om_\Theta$. \emph{Composition.} Using
$\intr\Om=\intr\Om_H\times\intr\Om_\Theta$ and Tonelli, for any $A$ and state $(\bu,\bth)$,
\[
  T\big((\bu,\bth),A\big)\ge\int_{\intr\Om_H}\delta_P\,\Leb(A_{\bu'}\cap\intr\Om_\Theta)\,\delta_H\,d\bu'
  =\delta_H\delta_P\,\Leb(A\cap\intr\Om).
\]
Set $\varepsilon:=\delta_H\delta_P\Leb(\Om)$ and $\nu(\cdot)=\Leb(\cdot\cap\intr\Om)/\Leb(\Om)$; then
$T\ge\varepsilon\nu$, and $A=\Om$ gives $\varepsilon\in(0,1]$ (as $\nu(\Om)=1$ by $\Leb(\partial\Om)=0$).
\end{proof}

\begin{theorem}[Uniform ergodicity / global convergence]
\label{thm:A}
Under Assumptions~\ref{ass:compact}--\ref{ass:cont}, $T$ has a unique invariant $\pi^\ast\in\cP(\Om)$
and, for every $\mu_0$, $\norm{\mu_0T^n-\pi^\ast}_\TV\le(1-\varepsilon)^n$ with
$\varepsilon=\delta_H\delta_P\Leb(\Om)\in(0,1]$. The rate is uniform over initializations.
\end{theorem}
\begin{proof}
$T$ is a kernel (Lemma~\ref{lem:measurable}); $\Om\subseteq\R^d$ has closed, hence Borel, diagonal;
minorization holds (Lemma~\ref{lem:minor}). Apply Theorem~\ref{thm:doeblin}.
\end{proof}

\begin{remark}[The TV constant is loose in high dimension]
\label{rem:scope}
From \eqref{eq:deltaH}, $\delta_H=(2\pi\sigma_H^2)^{-d_H/2}p_R\exp(-(M_H+M_H')^2/2\sigma_H^2)$, whose
prefactor and exponent are strongly dimension- and $\sigma$-dependent; so $\varepsilon$ is typically
minuscule and $(1-\varepsilon)^n$ certifies convergence but is a loose rate. Dimension-robust rates
come from the Wasserstein contraction below, on disjoint hypotheses.
\end{remark}

% =====================================================================
\section{Mixing rate: Wasserstein contraction (Theorem~\ref{thm:rate})}
\label{sec:rate}

\begin{assumption}[Conditional Lipschitz decoders --- (A-contraction)]
\label{ass:lip}
There exist $L_H,L_P<\infty$ with $\sup_{z}\mathrm{Lip}_{\bth\in\Om_\Theta}(g_\phi(\cdot,z))\le L_H$ and
$\sup_{z}\mathrm{Lip}_{\bu\in\Om_H}(g_\psi(\cdot,z))\le L_P$: each decoder mean is Lipschitz in its
\emph{conditioning} variable, uniformly in the latent. This is exactly the quantity condition-noise
injection regularizes by the noise-equals-Tikhonov identity \citep{bishop1995}; we posit it as an
assumption and support it empirically.
\end{assumption}

\begin{lemma}[Conditional $W_2$-Lipschitzness; the injection noise cancels]
\label{lem:condlip}
Under Assumptions~\ref{ass:compact},\ref{ass:lip}, $W_2(Q_H(\cdot\mid\bth),Q_H(\cdot\mid\bth'))\le
L_H\norm{\bth-\bth'}$ and $W_2(Q_P(\cdot\mid\bu),Q_P(\cdot\mid\bu'))\le L_P\norm{\bu-\bu'}$.
\end{lemma}
\begin{proof}
Couple by sharing latent $\bm z_H$ and injection noise $\bm\xi_H$:
$U=\mathcal B_H(g_\phi(\bth,\bm z_H)+\sigma_H\bm\xi_H)$, $U'=\mathcal B_H(g_\phi(\bth',\bm z_H)+\sigma_H\bm\xi_H)$.
As $\mathcal B_H$ is $1$-Lipschitz and $\sigma_H\bm\xi_H$ cancels,
$\norm{U-U'}\le\norm{g_\phi(\bth,\bm z_H)-g_\phi(\bth',\bm z_H)}\le L_H\norm{\bth-\bth'}$ pathwise, so
$W_2\le(\E\norm{U-U'}^2)^{1/2}\le L_H\norm{\bth-\bth'}$. Same for $Q_P$.
\end{proof}

\begin{theorem}[Sweep contraction]
\label{thm:rate}
Let $K(\bth,\cdot)=\int_{\Om_H}Q_P(\cdot\mid\bu)\,Q_H(d\bu\mid\bth)$ be the $\bth$-marginal sweep. Under
Assumptions~\ref{ass:compact}--\ref{ass:lip},
$W_2(K(\bth,\cdot),K(\bth',\cdot))\le\kappa\norm{\bth-\bth'}$ with $\kappa:=L_HL_P$, and the measure
action is a $\kappa$-contraction. If $\kappa<1$, $K$ has a unique invariant $\nu^\ast\in\cP(\Om_\Theta)$
with $W_2(\mu_0K^n,\nu^\ast)\le\kappa^n W_2(\mu_0,\nu^\ast)$, equal to the $\bth$-marginal of $\pi^\ast$.
\end{theorem}
\begin{proof}
Share all four noises. With $U,U'$ as in Lemma~\ref{lem:condlip} and
$\Theta'=\mathcal B_\Theta(g_\psi(U,\bm z_P)+\sigma_P\bm\xi_P)$,
$\tilde\Theta'=\mathcal B_\Theta(g_\psi(U',\bm z_P)+\sigma_P\bm\xi_P)$, the same cancellation gives
$\norm{\Theta'-\tilde\Theta'}\le L_P\norm{U-U'}\le\kappa\norm{\bth-\bth'}$ pathwise, so
$W_2(K(\bth,\cdot),K(\bth',\cdot))\le\kappa\norm{\bth-\bth'}$. This shared-noise coupling is an explicit
measurable map; gluing it against an optimal coupling of $\mu,\mu'$ gives
$W_2(\mu K,\mu'K)\le\kappa W_2(\mu,\mu')$. For $\kappa<1$, Fact~\ref{fact:std}(f) yields $\nu^\ast$ and
the rate; invariance of $\pi^\ast$ under $T$ identifies its $\bth$-marginal with $\nu^\ast$.
\end{proof}

\begin{remark}[The four noises; the collapse ``paradox'' dissolves]
\label{rem:noises}
$\kappa=L_HL_P$ is governed \emph{solely} by the condition-Lipschitz constants (regularized by
\citealp{bishop1995}); the injection noise $\sigma_H,\sigma_P$ and the latent \emph{cancel} in the
couplings and do not enter $\kappa$. Conversely, $\sigma>0$ is exactly what keeps each $Q$
nondegenerate (Lemma~\ref{lem:minor}) and sets the width of $\pi^\ast$; a deterministic map would
$W_2$-contract to a point mass---the mode collapse to avoid. Tightening the contraction and preventing
collapse thus act on \emph{disjoint} parts of the analysis: the apparent ``more contraction
$\Rightarrow$ more collapse'' tension was an artifact of conflating the condition-Lipschitz constant
with the injection scale.
\end{remark}

% =====================================================================
\section{Steering, consistency, and the certificate (Theorem B, $\einc$)}
\label{sec:steering}

\begin{assumption}[The domain contains the true support]
\label{ass:suptrue}
$\operatorname{supp}\Pi^\dagger\subseteq\Om$, where $\Pi^\dagger:=p(\bu,\bth\mid\xobs)$ is the true
joint posterior. Fix regular versions of the true conditionals $p^\dagger_H(\cdot\mid\bth)$,
$p^\dagger_P(\cdot\mid\bu)$ for every $\bth,\bu$.
\end{assumption}

\paragraph{The target.} Let $K^\dagger(\bth,\cdot)=\int p^\dagger_P(\cdot\mid\bu)\,p^\dagger_H(d\bu\mid\bth)$
be the true sweep. Since $p^\dagger_H,p^\dagger_P$ are conditionals of the one joint $\Pi^\dagger$
(hence compatible), Gibbs is exact: $\Pi^\dagger$ is $K^\dagger$-invariant and its $\bth$-marginal
$\nu^\dagger:=p(\bth\mid\xobs)$ is the estimand. Define the approximation defects
\begin{equation}\label{eq:defects}
  \varepsilon_H:=\sup_{\bth}W_2\big(Q_H(\cdot\mid\bth),p^\dagger_H(\cdot\mid\bth)\big),\qquad
  \varepsilon_P:=\sup_{\bu}W_2\big(Q_P(\cdot\mid\bu),p^\dagger_P(\cdot\mid\bu)\big).
\end{equation}

\begin{lemma}[Operator discrepancy]
\label{lem:opdisc}
Under Assumptions~\ref{ass:compact}--\ref{ass:lip},
$\sup_{\bth}W_2(K(\bth,\cdot),K^\dagger(\bth,\cdot))\le L_P\varepsilon_H+\varepsilon_P$.
\end{lemma}
\begin{proof}
Interpolate through $K'(\bth,\cdot)=\int Q_P(\cdot\mid\bu)\,p^\dagger_H(d\bu\mid\bth)$. \emph{Swap $H$:}
take an optimal coupling $(U,U^\dagger)$ of $Q_H(\cdot\mid\bth),p^\dagger_H(\cdot\mid\bth)$ and draw the
$P$-step by shared noise (Lemma~\ref{lem:condlip}), so
$\norm{\theta'-\theta'^\dagger}\le L_P\norm{U-U^\dagger}$, giving $W_2(K,K')\le L_P\varepsilon_H$.
\emph{Swap $P$:} with the same $\bu$, couple $Q_P(\cdot\mid\bu),p^\dagger_P(\cdot\mid\bu)$ by a
measurable family (Fact~\ref{fact:std}(e)) glued against $p^\dagger_H(\cdot\mid\bth)$, giving
$W_2(K',K^\dagger)\le\varepsilon_P$. Sum and take $\sup_{\bth}$.
\end{proof}

\begin{lemma}[Fixed-point perturbation]
\label{lem:perturb}
If $K$'s action is a $\kappa$-contraction ($\kappa\in[0,1)$) with invariant $\nu^\ast$ and $K'$ is any
kernel with invariant $\nu'$, then
$W_2(\nu^\ast,\nu')\le\frac{1}{1-\kappa}\sup_{\bth}W_2(K(\bth,\cdot),K'(\bth,\cdot))$.
\end{lemma}
\begin{proof}
$W_2(\nu^\ast,\nu')=W_2(\nu^\ast K,\nu'K')\le W_2(\nu^\ast K,\nu'K)+W_2(\nu'K,\nu'K')
\le\kappa W_2(\nu^\ast,\nu')+\sup_{\bth}W_2(K(\bth,\cdot),K'(\bth,\cdot))$, using
Fact~\ref{fact:std}(e) for the second term; rearrange.
\end{proof}

\begin{theorem}[Steering bound]
\label{thm:B}
Under Assumptions~\ref{ass:compact}--\ref{ass:lip} with $\kappa<1$,
\[
  W_2\big(\nu^\ast,\,p(\bth\mid\xobs)\big)\le\frac{L_P\varepsilon_H+\varepsilon_P}{1-\kappa}.
\]
In particular $\varepsilon_H=\varepsilon_P=0$ gives $\nu^\ast=p(\bth\mid\xobs)$ exactly (the
compatible/well-specified case).
\end{theorem}
\begin{proof}
$K$ is a $\kappa$-contraction with invariant $\nu^\ast$ (Theorem~\ref{thm:rate}), $K^\dagger$ has
invariant $\nu^\dagger$; apply Lemma~\ref{lem:perturb} and Lemma~\ref{lem:opdisc}.
\end{proof}

\paragraph{The computable self-consistency defect.} Let
$\mu^\ast(d\bu):=\int Q_H(d\bu\mid\bth)\,\nu^\ast(d\bth)$ and form, on $\Om_\Theta\times\Om_H$,
\begin{equation}\label{eq:einc}
  \Pi^{\rightarrow}(d\bth,d\bu):=\nu^\ast(d\bth)\,Q_H(d\bu\mid\bth),\quad
  \Pi^{\leftarrow}(d\bth,d\bu):=Q_P(d\bth\mid\bu)\,\mu^\ast(d\bu),\quad
  \einc:=W_2(\Pi^{\rightarrow},\Pi^{\leftarrow}).
\end{equation}
Both have $\bu$-marginal $\mu^\ast$; $\einc$ is the forward-vs-backward sweep inconsistency, estimable
from samples of the trained model alone.

\begin{proposition}[$\einc$ certifies compatibility]
\label{prop:inc}
Assume $\kappa<1$. Then $\einc=0$ iff $Q_H,Q_P$ are compatible (a.e.\ the conditionals of a common
joint), in which case $\Pi^\rightarrow=\Pi^\leftarrow$ and $\nu^\ast$ is its $\bth$-marginal.
\end{proposition}
\begin{proof}
$(\Leftarrow)$ A witnessing joint $\Pi'$ reproduces itself under one forward sweep, so
$\Pi'_\theta=\nu^\ast$ (unique) and $\Pi'=\nu^\ast\otimes Q_H=\Pi^\rightarrow=Q_P\otimes\mu^\ast=\Pi^\leftarrow$.
$(\Rightarrow)$ $\einc=0\Rightarrow\Pi^\rightarrow=\Pi^\leftarrow=:\Pi$, which has $\bu\mid\bth$-law
$Q_H$ and $\bth\mid\bu$-law $Q_P$, so they are compatible with witness $\Pi$.
\end{proof}

\begin{remark}[$\Pi^\leftarrow$ is the stationary joint]
\label{rem:pileft}
At stationarity the swept state has $\bu'\sim\mu^\ast$ and $\bth'\sim Q_P(\cdot\mid\bu')$, i.e.\
$\pi^\ast=\Pi^\leftarrow$ (up to the isometric coordinate swap
$\Om_H\times\Om_\Theta\cong\Om_\Theta\times\Om_H$). Hence $\einc=W_2(\Pi^\rightarrow,\pi^\ast)$: the
defect measures how far the forward pairing sits from the sampler's own stationary joint---a single-run
quantity.
\end{remark}

\begin{remark}[Why the defects are not summed --- no double counting]
\label{rem:nodouble}
Adding $\einc$ to Theorem~\ref{thm:B} would double count: because the true conditionals are compatible,
the forward and backward true joints coincide, so $\einc$ is \emph{dominated by} the approximation
defects rather than independent of them. The precise constant is Theorem~\ref{thm:incbound}.
\end{remark}

\paragraph{The constant.} We now discharge Remark~\ref{rem:nodouble} with an explicit constant. Beyond
the learned-model regularity in force, we need one hypothesis on the truth.

\begin{assumption}[Regularity of the true conditionals]
\label{ass:truereg}
The fixed versions of Assumption~\ref{ass:suptrue} are $W_2$-Lipschitz in their conditioning variable:
for some $L_H^\dagger,L_P^\dagger<\infty$,
$W_2(p^\dagger_H(\cdot\mid\bth),p^\dagger_H(\cdot\mid\bth'))\le L_H^\dagger\norm{\bth-\bth'}$ and
$W_2(p^\dagger_P(\cdot\mid\bu),p^\dagger_P(\cdot\mid\bu'))\le L_P^\dagger\norm{\bu-\bu'}$. This is the
true-posterior analogue of Assumption~\ref{ass:lip} and constrains only the \emph{truth}.
\end{assumption}

\begin{lemma}[Joint assembly bound]
\label{lem:assembly}
Let $\alpha,\alpha'\in\cP(\mathsf S)$, $M,M'$ Markov kernels from compact $\mathsf S$ to compact
$\Om'$, and $\Pi(ds,dy)=\alpha(ds)M(s,dy)$, $\Pi'(ds,dy)=\alpha'(ds)M'(s,dy)$ on $\mathsf S\times\Om'$
(Euclidean product metric). Then for \emph{every} $\gamma\in\Gamma(\alpha,\alpha')$,
\[
  W_2(\Pi,\Pi')^2\le\int\norm{s-s'}^2\gamma(ds,ds')+\int W_2\big(M(s,\cdot),M'(s',\cdot)\big)^2\gamma(ds,ds').
\]
\end{lemma}
\begin{proof}
Viewing $M,M'$ as kernels from $\mathsf S\times\mathsf S$, Fact~\ref{fact:std}(e) gives a measurable
family $(s,s')\mapsto\gamma^M_{s,s'}$ of optimal couplings of $M(s,\cdot),M'(s',\cdot)$. Then
$\Xi(ds,ds',dy,dy'):=\gamma(ds,ds')\gamma^M_{s,s'}(dy,dy')$ is a coupling of $\Pi,\Pi'$ (its $(s,y)$-
and $(s',y')$-marginals are $\Pi,\Pi'$), and evaluating $\int(\norm{s-s'}^2+\norm{y-y'}^2)\,d\Xi$ gives
the bound.
\end{proof}

\begin{theorem}[The self-consistency defect is dominated by the approximation defects]
\label{thm:incbound}
Under Assumptions~\ref{ass:compact}--\ref{ass:lip} with $\kappa<1$ and
Assumptions~\ref{ass:suptrue}--\ref{ass:truereg}, with $A:=2+L_P^\dagger$,
\[
  \einc\le C_H\varepsilon_H+C_P\varepsilon_P\le C(\varepsilon_H+\varepsilon_P),\quad
  C_H=A+\tfrac{A(1+L_H^\dagger)L_P}{1-\kappa},\ \
  C_P=1+\tfrac{A(1+L_H^\dagger)}{1-\kappa},\ \ C=\max(C_H,C_P).
\]
In particular $\einc=0$ whenever $\varepsilon_H=\varepsilon_P=0$, consistent with
Proposition~\ref{prop:inc}.
\end{theorem}
\begin{proof}
$W_2$ on products uses the Euclidean product metric; the isometric swap
$\Om_\Theta\times\Om_H\cong\Om_H\times\Om_\Theta$ (Remark~\ref{rem:pileft}) lets the three joints be
compared in one space. Set $D:=W_2(\nu^\ast,\nu^\dagger)$; Theorem~\ref{thm:B} gives
$D\le(L_P\varepsilon_H+\varepsilon_P)/(1-\kappa)$. Since $p^\dagger_H,p^\dagger_P$ are the conditionals
of the single joint $\Pi^\dagger$, its two disintegrations coincide,
\begin{equation}\label{eq:pivot}
  \Pi^\dagger=\nu^\dagger(d\bth)p^\dagger_H(d\bu\mid\bth)=\mu^\dagger(d\bu)p^\dagger_P(d\bth\mid\bu),
  \quad\mu^\dagger:=\text{$\bu$-marginal of }\Pi^\dagger .
\end{equation}

\emph{Step 1 (forward).} Apply Lemma~\ref{lem:assembly} with $\alpha=\nu^\ast,\alpha'=\nu^\dagger,
M=Q_H,M'=p^\dagger_H$ and an optimal $\gamma$ ($\int\norm{\bth-\bth'}^2\gamma=D^2$). The joints are
$\Pi^\rightarrow$ and, by \eqref{eq:pivot}, $\Pi^\dagger$. By the triangle inequality and
Assumption~\ref{ass:truereg}, $W_2(Q_H(\cdot\mid\bth),p^\dagger_H(\cdot\mid\bth'))\le\varepsilon_H+
L_H^\dagger\norm{\bth-\bth'}$; Minkowski in $L^2(\gamma)$ and $\sqrt{a^2+b^2}\le a+b$ give
\begin{equation}\label{eq:rH}
  W_2(\Pi^\rightarrow,\Pi^\dagger)\le\varepsilon_H+(1+L_H^\dagger)D=:r_H.
\end{equation}
\emph{Step 2 (backward).} The $\bu$-marginals of $\Pi^\rightarrow,\Pi^\dagger$ are $\mu^\ast,\mu^\dagger$;
marginalization is $1$-Lipschitz (Fact~\ref{fact:std}(d)), so $E:=W_2(\mu^\ast,\mu^\dagger)\le r_H$ by
\eqref{eq:rH}. Applying Lemma~\ref{lem:assembly} with $\alpha=\mu^\ast,\alpha'=\mu^\dagger,M=Q_P,
M'=p^\dagger_P$ (joints $\Pi^\leftarrow$ and $\Pi^\dagger$, by \eqref{eq:pivot}) and the analogous bound,
\begin{equation}\label{eq:rP}
  W_2(\Pi^\leftarrow,\Pi^\dagger)\le\varepsilon_P+(1+L_P^\dagger)E\le\varepsilon_P+(1+L_P^\dagger)r_H=:r_P.
\end{equation}
\emph{Step 3.} By the triangle inequality and $A=2+L_P^\dagger$,
$\einc\le r_H+r_P=A\,r_H+\varepsilon_P=A\varepsilon_H+\varepsilon_P+A(1+L_H^\dagger)D$. Substituting the
bound on $D$ and collecting coefficients gives $C_H,C_P$. If $\varepsilon_H=\varepsilon_P=0$ then $D=0$,
so $r_H=r_P=0$.
\end{proof}

\begin{remark}[Reading the constant; the one-sided guarantee]
\label{rem:incconst}
$C$ inflates only through $1/(1-\kappa)$ as $\kappa\uparrow1$, and grows with the true-conditional
roughness $L_H^\dagger,L_P^\dagger$ and the learned $L_P$; it does \emph{not} depend on $\sigma$ or
dimension directly (contrast $\varepsilon$, Remark~\ref{rem:scope}). Thus small $\varepsilon_H,
\varepsilon_P$ force small $\einc$, and---contrapositively---a large \emph{observed} $\einc$ certifies
that some conditional is inaccurate. The reverse fails: $\einc$ small does \emph{not} imply closeness to
truth (a model can be self-consistent yet wrong), so $\einc$ certifies \emph{compatibility}
(Proposition~\ref{prop:inc}), not \emph{accuracy}---the latter is Theorem~\ref{thm:B}. This is the
sense in which $\einc$ is a computable \emph{shadow} of the uncomputable $\varepsilon_H,\varepsilon_P$.
\end{remark}

\begin{remark}[Where the distortion defect lives]
\label{rem:edist}
As $p^\dagger_H,p^\dagger_P$ are the \emph{clean} true conditionals, $\varepsilon_H,\varepsilon_P$
silently include a distortion defect
$\edist=O(\sigma_H+\sigma_P)+(\E\norm{\bm Y-\mathcal B(\bm Y)}^2)^{1/2}$ (injection smoothing plus RMS
projection displacement). Referencing instead the $\sigma$-smoothed projected conditionals
$p^{\dagger,\sigma}_H,p^{\dagger,\sigma}_P$ isolates pure learning error and adds a single
$W_2(\nu^{\dagger,\sigma},\nu^\dagger)=O(\sigma)$ smoothing term, which has an irreducible floor since
$\sigma>0$ is required for Theorem~\ref{thm:A} and non-collapse; in the non-identifiable regime this
floor is small relative to the true ridge width, so the ridge is faithfully covered.
\end{remark}

% =====================================================================
\section{Positioning against approximate MCMC and dependency networks}
\label{sec:positioning}

The tools we use are standard, and we say so plainly: our steering bound (Theorem~\ref{thm:B}) has the
same shape---perturbation size over $1-\text{contraction}$---as classical \emph{perturbed-MCMC}
bounds. Novelty therefore lies not in the machinery but in the \emph{object} analyzed and in a
\emph{computable certificate}.

\paragraph{Approximate / perturbed MCMC.}
In perturbation theory for Markov chains \citep{mitrophanov2005,rudolf2018,johndrow2015,medina2016,negrea2021}
one fixes a single \emph{ideal} kernel $P$ whose invariant law is the known target $\pi$, and a single
\emph{approximate} kernel $\hat P$ actually simulated, bounding $\|\hat\pi-\pi\|$ by the operator
discrepancy over a contraction gap. Two structural features distinguish our setting. First, the object
perturbed is not one chain but a \emph{pair of separately learned conditionals} $Q_H,Q_P$ that
\emph{need not be the conditionals of any joint}; the very notion of (in)compatibility does not arise
with one chain approximating one target. Second, the perturbation ``size'' there, like our
$\varepsilon_H,\varepsilon_P$, is generally \emph{not computable} without the truth. Our contribution is
a certificate that is.

\paragraph{Dependency networks and conditional specification.}
Learning each conditional separately and running pseudo-Gibbs is the dependency-network recipe
\citep{heckerman2000}, and compatibility of a conditional specification is classical
\citep{arnold1999}; but those analyses are \emph{discrete} and \emph{asymptotic}. We give
\emph{continuous}-state, \emph{non-asymptotic}, quantitative $W_2$ statements and a computable
compatibility defect bounded to the approximation error.

\paragraph{Iterated random functions / Wasserstein ergodicity.}
The $W_2$ contraction is in the lineage of iterated random functions \citep{diaconis1999} and
Wasserstein/Harris ergodicity \citep{hairer2011}; we use it as a tool and credit it as such.

\paragraph{The certificate $\einc$, precisely.}
Table~\ref{tab:einc} states what $\einc$ does and does not guarantee. Its one-sided nature is the
honest boundary of the claim, and no analogue of such an in-run certificate---tying the uncomputable
perturbation size to an observable quantity with an explicit constant---exists in the perturbed-MCMC
literature.

\begin{table}[h]
\centering
\begin{tabular}{@{}lcl@{}}
\toprule
Statement & Holds? & Backing \\
\midrule
$\varepsilon_H,\varepsilon_P$ small $\Rightarrow \einc$ small & yes & Thm.~\ref{thm:incbound} \\
$\einc$ large $\Rightarrow$ some conditional inaccurate & yes & contrapositive of Thm.~\ref{thm:incbound} \\
$\einc=0 \Rightarrow$ sampling a genuine joint's marginal & yes & Prop.~\ref{prop:inc} \\
$\einc$ small $\Rightarrow$ close to the truth & \textbf{no} & (accuracy is Thm.~\ref{thm:B}, not $\einc$) \\
\bottomrule
\end{tabular}
\caption{What the computable certificate $\einc$ guarantees: a \emph{one-sided} shadow of the
approximation defects, never a substitute for the steering bound.}
\label{tab:einc}
\end{table}

% =====================================================================
\section{Numerical illustration}
\label{sec:numerics}

Three controlled experiments show the constants are real and the bounds meaningful. The first two are
linear-Gaussian, so the true posterior, $\kappa$, the defects $\varepsilon_H,\varepsilon_P$, and
$\einc$ are all available in closed form (Wasserstein distances between Gaussians via the Bures
formula); a Monte-Carlo run confirms them and demonstrates single-run estimability of $\einc$.

\paragraph{(E-a) Linear-Gaussian, compatible.}
A genuine bivariate Gaussian $(\bu,\bth)\sim\Normal(\bm0,\Sigma)$ with correlation $\rho$ supplies the
truth, and the learned conditionals are set equal to the true ones ($\varepsilon_H=\varepsilon_P=0$).
The condition-Lipschitz constants give $\kappa=L_HL_P=\rho^2$---the classical bivariate-Gaussian Gibbs
autocorrelation, an independent check---and choosing unequal marginal scales makes one individual
$L>1$ while $\kappa<1$, illustrating Remark~\ref{rem:noises}. \emph{Validates:} the geometric rate
$\kappa^n$ (Theorem~\ref{thm:rate}) via the exact AR(1) $\bth$-chain, and exact recovery
$\nu^\ast=\nu^\dagger$ (Theorem~\ref{thm:B}).

\paragraph{(E-b) Controlled incompatibility.}
Perturbing one conditional's slope by $\delta$ (keeping $\varepsilon_H=0$, dialing
$\varepsilon_P=\delta U$ on the box of radius $U$) makes the pair incompatible while keeping the chain
Gaussian. \emph{Validates:} that the steering bound $\varepsilon_P/(1-\kappa)$ upper-bounds the true
$W_2(\nu^\ast,\nu^\dagger)$ for all $\delta$ (Theorem~\ref{thm:B}); that the closed-form $\einc(\delta)$
stays below $C(\delta)\,\varepsilon_P(\delta)$ and co-vanishes with, and tracks, $\varepsilon_P$
(Theorem~\ref{thm:incbound}); and that $\einc$ estimated from a single run matches the closed form.

\paragraph{(E-c) Non-identifiable toy.}
A product-degenerate forward map $\bu=a\cdot b$ (parameters $\bth=(a,b)$) yields a hyperbolic-ridge
posterior---the pure-mathematics reduction of the application's $S_I\!\cdot\!\sigma$ fiber. With learned
$=$ true conditionals, the experiment isolates non-identifiability. \emph{Validates:} collapse-free
support recovery (the stochastic dual sweep covers the ridge to width $O(\sigma_H+\sigma_P)$,
Remark~\ref{rem:edist}) against a deterministic ping-pong that collapses to a point;
$\einc\approx0$ certifies compatibility on the ridge (Proposition~\ref{prop:inc}); and a $2$-D grid
supplies the reference posterior. Full design in the companion notes.

% =====================================================================
\section{Discussion and limitations}
\label{sec:discussion}

\paragraph{The contraction hypothesis.} Existence and uniform ergodicity (Theorem~\ref{thm:A}) are
proved from what the algorithm controls---compactness and $\sigma>0$---with \emph{no} contraction
assumed. The rate and the steering constant, however, rest on the condition-Lipschitz bound
$\kappa=L_HL_P<1$ (Assumption~\ref{ass:lip}), which training \emph{encourages} but does not
\emph{certify}: condition-noise injection penalizes $\E\|\nabla_{\text{cond}}g\|_F^2$
\citep{bishop1995}, promoting small $L_H,L_P$, and one can treat $\kappa$ as a free constant and
support $\kappa<1$ empirically through autocorrelation decay and the monitored $\einc$. We regard this
as an honest assumption rather than a theorem, and phrase every rate/steering statement conditionally
on it.

\paragraph{The $\sigma$-floor.} The injection noise $\sigma_H,\sigma_P>0$ is doubly load-bearing: it
supplies the minorization of Theorem~\ref{thm:A} and prevents mode collapse, yet it also smooths each
conditional, contributing an irreducible $O(\sigma)$ term to the distortion defect
(Remark~\ref{rem:edist}). This floor cannot be removed without losing the guarantees; the reconciling
point is that the target is a \emph{distribution}, not a point, and in the non-identifiable regime the
floor is small relative to the width of the true posterior ridge, so the ridge is faithfully covered
(experiment E-c).

\paragraph{The certificate is one-sided.} $\einc$ certifies \emph{compatibility}, not \emph{accuracy}:
$\einc=0$ means the sampler draws from a genuine joint's marginal (Proposition~\ref{prop:inc}), and a
large observed $\einc$ is a rigorous witness of an inaccurate conditional (contrapositive of
Theorem~\ref{thm:incbound}), but a small $\einc$ does not imply closeness to the truth. Accuracy is the
separate content of Theorem~\ref{thm:B}, whose defects $\varepsilon_H,\varepsilon_P$ require the
unknown truth. Reading $\einc$ as an accuracy certificate would be a mistake; it is a computable shadow.

\paragraph{Extensions.} Several directions are natural: a joint $(\bu,\bth)$ contraction under a
weighted product metric (constant still governed by $L_H,L_P$); non-compact domains via a
Lyapunov/drift condition in place of the boundary projection; injection noise beyond Gaussian
(any nondegenerate, positive-density family suffices for the minorization); quantifying
$\varepsilon_H,\varepsilon_P$ from finite data and capacity, linking to simulation-based inference
\citep{papamakarios2016,cranmer2020}; and relaxing the true-conditional regularity
(Assumption~\ref{ass:truereg}) to local or integrated Lipschitz control. The bounds here are
worst-case (suprema over the conditioning variable) and the constants can be loose; sharpening them,
and characterizing when $\einc$ is not merely an upper but also a lower proxy for the approximation
error, are open.

% =====================================================================
\begin{thebibliography}{99}

\bibitem[Arnold et al.(1999)]{arnold1999}
B.~C. Arnold, E.~Castillo, and J.~M. Sarabia.
\emph{Conditional Specification of Statistical Models}. Springer, 1999.

\bibitem[Bishop(1995)]{bishop1995}
C.~M. Bishop.
Training with noise is equivalent to Tikhonov regularization.
\emph{Neural Computation}, 7(1):108--116, 1995.

\bibitem[Cranmer et al.(2020)]{cranmer2020}
K.~Cranmer, J.~Brehmer, and G.~Louppe.
The frontier of simulation-based inference.
\emph{PNAS}, 117(48):30055--30062, 2020.

\bibitem[Diaconis and Freedman(1999)]{diaconis1999}
P.~Diaconis and D.~Freedman.
Iterated random functions.
\emph{SIAM Review}, 41(1):45--76, 1999.

\bibitem[Folland(1999)]{folland}
G.~B. Folland.
\emph{Real Analysis: Modern Techniques and Their Applications}, 2nd ed. Wiley, 1999.

\bibitem[Hairer et al.(2011)]{hairer2011}
M.~Hairer, J.~C. Mattingly, and M.~Scheutzow.
Asymptotic coupling and a general form of Harris' theorem for Markov processes.
\emph{Probability Theory and Related Fields}, 149(1--2):223--259, 2011.

\bibitem[Heckerman et al.(2000)]{heckerman2000}
D.~Heckerman, D.~M. Chickering, C.~Meek, R.~Rounthwaite, and C.~Kadie.
Dependency networks for inference, collaborative filtering, and data visualization.
\emph{Journal of Machine Learning Research}, 1:49--75, 2000.

\bibitem[Hiriart-Urruty and Lemar\'echal(2001)]{hiriart}
J.-B. Hiriart-Urruty and C.~Lemar\'echal.
\emph{Fundamentals of Convex Analysis}. Springer, 2001.

\bibitem[Johndrow et al.(2015)]{johndrow2015}
J.~E. Johndrow, J.~C. Mattingly, S.~Mukherjee, and D.~Dunson.
Approximations of Markov chains and high-dimensional Bayesian inference.
\emph{arXiv:1508.03387}, 2015.

\bibitem[Kallenberg(2002)]{kallenberg}
O.~Kallenberg.
\emph{Foundations of Modern Probability}, 2nd ed. Springer, 2002.

\bibitem[Medina-Aguayo et al.(2016)]{medina2016}
F.~J. Medina-Aguayo, A.~Lee, and G.~O. Roberts.
Stability of noisy Metropolis--Hastings.
\emph{Statistics and Computing}, 26:1187--1211, 2016.

\bibitem[Meyn and Tweedie(2009)]{meyntweedie}
S.~Meyn and R.~L. Tweedie.
\emph{Markov Chains and Stochastic Stability}, 2nd ed. Cambridge Univ.\ Press, 2009.

\bibitem[Mitrophanov(2005)]{mitrophanov2005}
A.~Yu. Mitrophanov.
Sensitivity and convergence of uniformly ergodic Markov chains.
\emph{Journal of Applied Probability}, 42(4):1003--1014, 2005.

\bibitem[Negrea and Rosenthal(2021)]{negrea2021}
J.~Negrea and J.~S. Rosenthal.
Approximations of geometrically ergodic reversible Markov chains.
\emph{Advances in Applied Probability}, 53(4):981--1022, 2021.

\bibitem[Papamakarios and Murray(2016)]{papamakarios2016}
G.~Papamakarios and I.~Murray.
Fast $\varepsilon$-free inference of simulation models with Bayesian conditional density estimation.
\emph{NeurIPS}, 2016.

\bibitem[Rudolf and Schweizer(2018)]{rudolf2018}
D.~Rudolf and N.~Schweizer.
Perturbation theory for Markov chains via Wasserstein distance.
\emph{Bernoulli}, 24(4A):2610--2639, 2018.

\bibitem[Villani(2009)]{villani}
C.~Villani.
\emph{Optimal Transport: Old and New}. Springer, 2009.

\end{thebibliography}

\end{document}
